\documentclass[11pt]{article}
\usepackage[T1]{fontenc}
\usepackage{lmodern}
\usepackage[margin=1.1in]{geometry}
\usepackage{amsmath, amssymb, amsthm}
\usepackage{booktabs}
\usepackage{xcolor}
\definecolor{linknavy}{RGB}{25, 55, 125}
\usepackage[colorlinks=true, linkcolor=linknavy, citecolor=linknavy, urlcolor=linknavy]{hyperref}
\hypersetup{
  pdftitle={Which counterexamples to the Petersen coloring conjecture are colorable only by themselves?},
  pdfauthor={Felipe Santibanez-Leal},
  pdfkeywords={Petersen coloring conjecture, H-coloring, cubic graphs, snarks, splitting lemma, SAT, DRAT}
}
\usepackage{orcidlink}
\usepackage{fancyhdr}

\newtheorem{theorem}{Theorem}[section]
\newtheorem{proposition}[theorem]{Proposition}
\newtheorem{lemma}[theorem]{Lemma}
\newtheorem{corollary}[theorem]{Corollary}
\theoremstyle{definition}
\newtheorem{definition}[theorem]{Definition}
\theoremstyle{remark}
\newtheorem{remark}[theorem]{Remark}

\newcommand{\PG}{P}
\newcommand{\Hthree}{\mathcal{H}_3}
\newcommand{\Ga}{G_{52}}
\newcommand{\Gb}{G'_{52}}
\newcommand{\Gsix}{G_{68}}
\newcommand{\Gmain}{G_{112}}
\newcommand{\Gsym}{H_{112}}

\newcommand{\doctype}{Preprint}
\newcommand{\docversion}{v0.01}
\newcommand{\docdate}{18 September 2026}
\newcommand{\versiondoi}{10.5281/zenodo.00000000}
\newcommand{\conceptdoi}{10.5281/zenodo.00000000}
\newcommand{\docrepo}{github.com/fsantibanezleal/CAOS\_RESEARCH}
\newcommand{\headerblock}{%
\begin{center}
\fbox{\parbox{0.94\textwidth}{\small\raggedright
\textbf{\doctype} \quad$\cdot$\quad \docversion, \docdate \quad$\cdot$\quad Licence CC BY 4.0\\[2pt]
CAOS Research open-problems programme \quad$\cdot$\quad problem \texttt{petersen-coloring}\\[2pt]
Version DOI \href{https://doi.org/\versiondoi}{\texttt{\versiondoi}}\\[2pt]
Concept DOI (always latest) \href{https://doi.org/\conceptdoi}{\texttt{\conceptdoi}}\\[2pt]
Not peer reviewed. Source code, data and computational records: \href{https://\docrepo}{\texttt{\docrepo}}
}}
\end{center}
\vspace{4pt}}

\title{Which Counterexamples to the Petersen Coloring Conjecture\\
Are Colorable Only by Themselves?}
\author{Felipe Santib\'a\~nez-Leal\,\orcidlink{0000-0002-0150-3246}\\
\small CAOS Research program, \texttt{github.com/fsantibanezleal/CAOS\_RESEARCH}}
\date{\docdate}

\pagestyle{fancy}
\fancyhf{}
\fancyhead[L]{\small CAOS Research \doctype}
\fancyhead[R]{\small Colorable Only by Itself \textperiodcentered\ \docversion}
\fancyfoot[C]{\small \thepage}
\renewcommand{\headrulewidth}{0.4pt}
\fancypagestyle{plain}{\fancyhf{}\fancyfoot[C]{\small \thepage}\renewcommand{\headrulewidth}{0pt}}

\begin{document}
\maketitle
\headerblock

\begin{abstract}
For cubic graphs $G$ and $H$, an $H$-coloring of $G$ is a map from the edges of $G$ to the edges
of $H$ that sends every vertex star of $G$ bijectively onto a vertex star of $H$; a Petersen
coloring is the case where $H$ is the Petersen graph. Ma, Mattiolo, Steffen and Wolf showed that
there is a unique minimal set $\mathcal{H}_3$ of connected bridgeless cubic graphs that color
every bridgeless cubic graph, and that a graph belongs to it exactly when no bridgeless cubic
graph of smaller order colors it. Since the Petersen coloring conjecture was refuted in 2026,
$\mathcal{H}_3$ is infinite and contains every smallest counterexample. Goedgebeur, Jooken,
M\'a\v{c}ajov\'a, Mattiolo, Mazzuoccolo and Ulyanov asked whether their two counterexamples on 52
vertices, the smallest known, are colorable only by themselves. We prove two lemmas on the vertex
map of an $H$-coloring: all its fibers have the same parity, and target vertices that are not
used can be reduced to at most one by the splitting lemma. They leave three kinds of colorings,
in none of which the target graph has a part unconstrained by $G$, and they turn the question
into one propositional formula per target order, with the target graph as part of the unknown.
For both 52-vertex counterexamples, every even target order from 40 to 50 is refuted with a DRAT
proof checked by \texttt{drat-trim}; several smaller orders are refuted as well. Together with the
result of Goedgebeur et al.\ that every bridgeless cubic graph on at most 38 vertices has a
Petersen coloring, this shows that both graphs belong to $\mathcal{H}_3$: they are colorable
only by themselves. Without the two lemmas the same search decided no target order of the first
graph within comparable time.
\end{abstract}


\section{Introduction}

\subsection{Colorings by a cubic graph}

Graphs are finite; they may have parallel edges but no loops. For a vertex $v$ of a graph $G$,
$\partial_G(v)$ is the set of edges at $v$. Let $G$ and $H$ be cubic graphs. An
\emph{$H$-coloring} of $G$ is a map $f : E(G) \to E(H)$ such that for every vertex $v$ of $G$
there is a vertex $u$ of $H$ with $f(\partial_G(v)) = \partial_H(u)$; since both stars have three
edges, $f$ is injective on every star. If such a map exists we say that $H$ \emph{colors} $G$. The
relation is transitive. A \emph{Petersen coloring} is a $\PG$-coloring, where $\PG$ is the
Petersen graph, and a cubic graph is 3-edge-colorable if and only if it is colored by the graph on
two vertices with three parallel edges.

Jaeger~\cite{Jaeger1988} conjectured that every bridgeless cubic graph has a Petersen coloring.
The conjecture implies the Berge--Fulkerson conjecture and the 5-cycle double cover conjecture.
It was refuted in August 2026: Putman~\cite{Putman2026} gave two counterexamples $\Gmain$ and
$\Gsym$ on 112 vertices with computer-checked certificates, Jooken gave a proof that can be
checked by hand, and Goedgebeur, Jooken, M\'a\v{c}ajov\'a, Mattiolo, Mazzuoccolo and
Ulyanov~\cite{GJMMMU2026} gave two counterexamples $\Ga$ and $\Gb$ on 52 vertices with a purely
theoretical proof, cyclically 4-edge-connected counterexamples of every even order at least 60,
and the bounds $40 \le n_{\min} \le 52$ for the order $n_{\min}$ of a smallest counterexample. A
counterexample $\Gsix$ on 68 vertices is also available from the House of
Graphs~\cite{HoG}.

\subsection{The set \texorpdfstring{$\Hthree$}{H3}}

Ma, Mattiolo, Steffen and Wolf~\cite{MMSW2025} proved that there is a unique inclusion-wise
minimal set $\Hthree$ of connected bridgeless cubic graphs such that every connected bridgeless
cubic graph is colored by a member of $\Hthree$, and they characterized its members.

\begin{theorem}[\cite{MMSW2025}, Theorem~3.7, case $r = 3$]\label{thm:mmsw}
For a connected bridgeless cubic graph $G$ the following are equivalent: $G \in \Hthree$; the only
connected bridgeless cubic graph that colors $G$ is $G$ itself; no bridgeless cubic graph of
smaller order colors $G$.
\end{theorem}

The Petersen coloring conjecture was the statement $\Hthree = \{\PG\}$. By~\cite{MMSW2025},
$\Hthree$ is either $\{\PG\}$ or infinite, so it is now known to be infinite, and its members other
than $\PG$ are all counterexamples: if $H$ colors a counterexample $G$ then $H$ is a
counterexample, because a Petersen coloring of $H$ composed with the $H$-coloring of $G$ would be
a Petersen coloring of $G$. In particular a smallest counterexample belongs to $\Hthree$.
Goedgebeur et al.\ ask~\cite[Section~5.4]{GJMMMU2026} whether their counterexamples of order 52
are colorable only by themselves, and observe that the question is tied to whether these graphs
are smallest counterexamples.

The question has two informative answers. A bridgeless cubic graph of order less than 52 that
colors $\Ga$ or $\Gb$ would be a counterexample below the known upper bound. If no such graph
exists, the 52-vertex graph is a member of $\Hthree$.

\subsection{Results}

We first prove two lemmas on the fibers of the vertex map of an $H$-coloring
(Section~\ref{sec:fibers}). All fibers have the same parity (Lemma~\ref{lem:parity}), and target
vertices that are not used by the coloring can be reduced to at most one by the splitting lemma
(Lemma~\ref{lem:unused}). Together they reduce the question ``is $G$ colored by some bridgeless
cubic graph of smaller order'' to finitely many propositional formulas, one per target order,
whose unknown target graph has no part left unconstrained by $G$
(Corollary~\ref{cor:modes}). Section~\ref{sec:encoding} describes the formulas and how each answer
is certified. Section~\ref{sec:results} reports the computations.

\section{Fibers of an \texorpdfstring{$H$}{H}-coloring}\label{sec:fibers}

Throughout this section $G$ is a connected cubic graph of order $n$, $H$ is a connected cubic
graph, $f$ is an $H$-coloring of $G$, and $\varphi : V(G) \to V(H)$ is a \emph{vertex map} of $f$:
a map with $f(\partial_G(v)) = \partial_H(\varphi(v))$ for every $v$. A vertex map exists by
definition, and it is unique unless $H$ is the graph with two vertices and three parallel edges.
For $x \in V(H)$ the \emph{fiber} of $x$ is $\varphi^{-1}(x)$ and $n_x = |\varphi^{-1}(x)|$; the
vertex $x$ is \emph{used} if $n_x > 0$. We write $U$ for the set of used vertices and
$W = V(H) \setminus U$.

\begin{lemma}[fiber parity]\label{lem:parity}
For every edge $e = xy$ of $H$, the set $f^{-1}(e)$ is a perfect matching of the subgraph of $G$
induced on $\varphi^{-1}(x) \cup \varphi^{-1}(y)$. Consequently $n_x + n_y$ is even for every edge
$xy$ of $H$, and all the numbers $n_x$, $x \in V(H)$, have the same parity.
\end{lemma}

\begin{proof}
A vertex $v$ of $G$ is incident with an edge of $f^{-1}(e)$ if and only if
$e \in f(\partial_G(v)) = \partial_H(\varphi(v))$, that is, if and only if
$\varphi(v) \in \{x, y\}$. In that case it is incident with exactly one such edge, because $f$ is
injective on $\partial_G(v)$. Hence $f^{-1}(e)$ is a matching that covers exactly
$\varphi^{-1}(x) \cup \varphi^{-1}(y)$, a set of $n_x + n_y$ vertices. The last statement follows
because $H$ is connected.
\end{proof}

So an $H$-coloring is of one of two kinds. If the fibers are odd, every vertex of $H$ is used and
$|V(H)| \le n$. If the fibers are even, at most $n/2$ vertices of $H$ are used. Both kinds occur:
an isomorphism has fibers of size~1, and a 3-edge-coloring of a graph of order $n \equiv 0
\pmod 4$ has a vertex map with two fibers of size $n/2$.

The second lemma concerns the vertices of $H$ that are not used. We use the splitting lemma of
Fleischner~\cite{Fleischner1992} in the following form. Let $v$ be a vertex of a graph $K$ and
$e_1 = vv_1$, $e_2 = vv_2$ two edges at $v$. \emph{Splitting $e_1$ and $e_2$ off $v$} means adding
a new vertex $v^*$ and replacing $e_1, e_2$ by edges $v^*v_1$, $v^*v_2$.

\begin{lemma}[\cite{KaiserKLW2007}, Lemma~1, a consequence of~\cite{Fleischner1992}]\label{lem:split}
Let $v$ be a vertex of degree at least~4 in a bridgeless graph $K$. There exist edges $e_1, e_2$
incident with $v$ such that the graph obtained by splitting $e_1$ and $e_2$ off $v$ is bridgeless.
\end{lemma}

\begin{lemma}[unused vertices]\label{lem:unused}
Let $H$ be bridgeless and suppose $W \ne \emptyset$. Let $m$ be the number of edges between $U$
and $W$. Then there is a connected bridgeless cubic graph $H'$ with vertex set $U \cup W'$, where
$W' = \emptyset$ if $m$ is even and $|W'| = 1$ if $m$ is odd, and an $H'$-coloring of $G$ with
the same vertex map $\varphi$. In particular $|V(H')| \le |V(H)|$.
\end{lemma}

\begin{proof}
Every edge of $H$ with an end in $U$ lies in the image of $f$, since the star of a used vertex is
the image of a star of $G$; no edge with both ends in $W$ does. Identify $W$ to a single vertex
$w^*$ and delete the loops that arise. Every edge cut of the resulting graph $H_1$ is an edge cut
of $H$, so $H_1$ is bridgeless; $w^*$ has degree $m$, and $m \ge 2$ because $H$ is connected and
bridgeless. Every other vertex of $H_1$ is a used vertex, of degree~3, with its star unchanged.

While $w^*$ has degree at least~4, apply Lemma~\ref{lem:split} to split two edges $aw^*$, $bw^*$
off $w^*$ so that the graph stays bridgeless. Then $a \ne b$: otherwise the new vertex would be
joined to the cubic vertex $a$ by two parallel edges, and the third edge at $a$ would be a bridge.
Suppress the new vertex, that is, replace its two edges by a single edge $ab$. This keeps the graph
loopless and bridgeless and lowers the degree of $w^*$ by~2. When the degree of $w^*$ has dropped
to~3, stop. When it has dropped to~2, suppress $w^*$ in the same way; its two neighbours are
distinct by the same argument. Let $H'$ be the final graph. It is cubic, loopless and bridgeless.

Let $\mu$ map every edge of $H$ that has an end in $U$ to the edge of $H'$ that contains it. Two
edges merged in one step have their ends other than $w^*$ at distinct vertices, so $\mu$ is
injective on $\partial_H(u)$ and $\partial_{H'}(u) = \mu(\partial_H(u))$ for every $u \in U$. Hence
$\mu \circ f$ is an $H'$-coloring of $G$ with vertex map $\varphi$.

Finally $H'$ is connected. If $uv$ is an edge of $G$ with $\varphi(u) \ne \varphi(v)$, then
$f(uv)$ joins $\varphi(u)$ to $\varphi(v)$; it has both ends in $U$ and is therefore an edge of
$H'$. As $G$ is connected, $U$ induces a connected subgraph of $H'$, and the vertex $w^*$, if it
remains, is joined to $U$.
\end{proof}

\begin{corollary}\label{cor:modes}
A connected cubic graph $G$ of order $n$ is colored by a connected bridgeless cubic graph of order
less than $n$ if and only if, for some even $k < n$, it has an $H$-coloring by a connected
bridgeless cubic graph $H$ of order $k$ with a vertex map of one of the following kinds:
\begin{enumerate}
\item[\textup{(O)}] every vertex of $H$ is used and every fiber is odd;
\item[\textup{(E0)}] every vertex of $H$ is used and every fiber is even, so that $k \le n/2$;
\item[\textup{(E1)}] exactly one vertex of $H$ is not used and every other fiber is even and
nonempty, so that $k - 1 \le n/2$.
\end{enumerate}
\end{corollary}

\begin{proof}
Let $H$ of order less than $n$ color $G$. If every vertex of $H$ is used, Lemma~\ref{lem:parity}
gives (O) or (E0). Otherwise the fibers are even by Lemma~\ref{lem:parity}, since an unused vertex
has an empty fiber, and Lemma~\ref{lem:unused} gives a graph $H'$ of order at most $|V(H)|$ with
the same vertex map and at most one unused vertex, which is (E0) or (E1).
\end{proof}

The kinds (O), (E0), (E1) leave no part of $H$ unconstrained by $G$: every edge of $H$ is the image
of an edge of $G$. For target orders close to $n$ only kind (O) is possible, and the fibers are then
almost all singletons. The extreme case of kind (O) can be settled by hand.

\begin{proposition}\label{prop:nminus2}
Let $G$ be a connected cubic graph of order $n \ge 8$ without triangles. Then $G$ has no
$H$-coloring by a connected cubic graph $H$ of order $n - 2$ in which every vertex of $H$ is used.
\end{proposition}

\begin{proof}
Since $n - 2 > n/2$ and every vertex of $H$ is used, the fibers are odd by
Lemma~\ref{lem:parity}; hence one fiber $\{u_1, u_2, u_3\} = \varphi^{-1}(x)$ has size~3 and all
others have size~1. For each of the three edges $e = xy$ at $x$, the set $f^{-1}(e)$ is a perfect matching on
$\{u_1, u_2, u_3\}$ together with the single vertex of the fiber of $y$
(Lemma~\ref{lem:parity}), so it contains an edge of $G$ joining two of $u_1, u_2, u_3$. The three
sets $f^{-1}(e)$ are pairwise disjoint and $G$ is simple, so $u_1u_2u_3$ is a triangle of $G$.
\end{proof}

\section{Propositional encoding and certification}\label{sec:encoding}

For a fixed graph $G$ and a fixed even $k$ we build one formula in conjunctive normal form that is
satisfiable if and only if $G$ has an $H$-coloring of kind (O), (E0) or (E1) by some loopless
cubic graph $H$ on $k$ vertices; connectedness and bridgelessness of $H$ are handled separately
below. The target graph is part of the unknown.

\begin{itemize}
\item \emph{Target.} Every target vertex $i < k$ carries three slots $(i, 0), (i, 1), (i, 2)$. A
perfect matching on the $3k$ slots that never pairs two slots of the same vertex is the edge set of
a loopless cubic graph on $k$ vertices, and every such graph arises. One variable per admissible
slot pair.
\item \emph{Map.} Variables $x_{v,i}$ state $\varphi(v) = i$, exactly one $i$ per vertex $v$. At
every $v$ a bijection between the three edges at $v$ and the three slots of $\varphi(v)$ records
$f$ on the star of $v$.
\item \emph{Consistency.} For every edge $uv$ of $G$, the slot assigned to $uv$ at $u$ and the slot
assigned to it at $v$ are either the same slot or a matched pair: both describe the edge $f(uv)$.
\item \emph{Symmetry breaking.} Target vertices are numbered in the order in which the vertices
$0, 1, \dots, n-1$ of $G$ first use them, and the first vertex of each fiber fixes the order of the
three slots.
\item \emph{Fibers.} One variable $q$ states that the fibers are odd; for every target vertex an
exclusive-or chain forces the parity of its fiber to equal $q$ (Lemma~\ref{lem:parity}). The last
target vertex but one is used (Lemma~\ref{lem:unused}), and $q$ is set to true when
$2(k - 1) > n$.
\end{itemize}

A model is decoded into a pair $(H, f)$. If $H$ is disconnected or has a bridge, a clause is added
that requires a further matched slot pair across the offending vertex partition, which every
connected bridgeless target satisfies, and the formula is solved again. The loop ends with a model
whose target is connected and bridgeless, or with an unsatisfiable formula.

\emph{Certification.} A satisfiable answer is checked from the definition by a separate program
that reads only $G$ and the decoded pair $(H, f)$: $H$ is cubic, loopless, connected and
bridgeless, and $f$ maps every star of $G$ onto a star of $H$. An unsatisfiable answer is
certified on the final formula (the base formula together with every added clause): it is refuted
by CaDiCaL~\cite{CaDiCaL} with a DRAT proof, and the proof is checked by
\texttt{drat-trim}~\cite{DRATtrim}. The incremental solver that drives the loop is not part of the
trusted base.

\emph{Controls.} The formulas with and without the fiber constraints were compared on every even
order below the order of the graph for $K_4$, the prism, the Petersen graph and the flower snarks
$J_3$ and $J_5$ (21 instances). The answers agree in every instance. In particular the Petersen
graph is colored by no bridgeless cubic graph of smaller order, $J_5$ is colored by bridgeless
cubic graphs of orders 10 and 12 and of no other order below 20, and $J_3$ is colored by one of
order 10 and by none of order 2, 4, 6 or 8.

\section{Results}\label{sec:results}

The graphs are $\Ga$ and $\Gb$ (House of Graphs entries 57244 and 57278), both cubic of girth~5
and cyclically 4-edge-connected. Their non-colorability by the Petersen graph was confirmed
independently of the published proofs: the edge-map formula with target $\PG$ is refuted with
checked proofs for both graphs, and so is the formula for normal 5-edge-colorings.

\begin{theorem}[machine-verified]\label{thm:orders}
For every even $k$ listed in Table~\ref{tab:orders}, no loopless cubic graph on $k$ vertices colors
the corresponding graph with a vertex map of kind \textup{(O)}, \textup{(E0)} or \textup{(E1)}.
The list contains every even $k$ with $40 \le k \le 50$ for both $\Ga$ and $\Gb$%%EXTRA_ORDERS%%.
\end{theorem}

Goedgebeur et al.\ established by an exhaustive computation that every bridgeless cubic graph on
at most 38 vertices has a Petersen coloring~\cite[Observation~9]{GJMMMU2026}. The statement
extends to graphs with parallel edges.

\begin{lemma}\label{lem:digon}
If a connected bridgeless cubic graph with parallel edges has no Petersen coloring, then a
connected bridgeless cubic graph of smaller order has none either.
\end{lemma}

\begin{proof}
Let $x, y$ be joined by two parallel edges; they are not joined by three, since the graph on two
vertices is Petersen colorable. Let $xx'$ and $yy'$ be the remaining edges at $x$ and $y$. Then
$x' \ne y'$, because otherwise the third edge at $x'$ would be a bridge. Delete $x$ and $y$ and
add the edge $x'y'$. Every edge cut of the new graph corresponds to an edge cut of the same size
of the old one, so the new graph is connected, bridgeless and cubic. If it had a Petersen
coloring $\sigma$ with $\sigma(x'y') = m$, then labelling $xx'$ and $yy'$ by $m$ and the two
parallel edges by the two other edges at one end of $m$ would give a Petersen coloring of the
old graph.
\end{proof}

\begin{theorem}\label{thm:h3}
The graphs $\Ga$ and $\Gb$ are colorable only by themselves: both belong to $\Hthree$.
\end{theorem}

\begin{proof}
Let $G$ be one of the two graphs and suppose that some connected bridgeless cubic graph of order
less than 52 colors $G$. By Corollary~\ref{cor:modes} there is one, say $H$, of even order
$k < 52$ with a vertex map of kind (O), (E0) or (E1). Since $H$ colors a graph without a Petersen
coloring, $H$ has no Petersen coloring. By Lemma~\ref{lem:digon}, applied repeatedly, there is a
simple connected bridgeless cubic graph of order at most $k$ without a Petersen coloring, so
$k \ge 40$ by~\cite[Observation~9]{GJMMMU2026}. This contradicts Theorem~\ref{thm:orders}.
Theorem~\ref{thm:mmsw} gives the conclusion.
\end{proof}

Theorem~\ref{thm:orders} does not depend on the computation of~\cite{GJMMMU2026}; only the orders
below 40 in Theorem~\ref{thm:h3} do. The orders below 40 that are listed in
Table~\ref{tab:orders} were refuted directly, and the remaining ones were not decided within six
hours each.


\begin{table}[h]\centering\small
\begin{tabular}{lrrrrrr}
\toprule
graph & $k$ & variables & clauses & solve (s) & proof (bytes) & check (s)\\
\midrule
$\Ga$ & 50 & 67{,}092 & 1{,}916{,}592 & 46 & 17{,}999{,}280 & 64 \\
$\Ga$ & 48 & 63{,}129 & 1{,}769{,}986 & 64 & 28{,}953{,}154 & 127 \\
$\Ga$ & 46 & 59{,}274 & 1{,}629{,}212 & 208 & 44{,}331{,}332 & 250 \\
$\Ga$ & 44 & 55{,}527 & 1{,}494{,}270 & 293 & 49{,}206{,}463 & 270 \\
$\Ga$ & 42 & 51{,}888 & 1{,}365{,}160 & 347 & 60{,}624{,}579 & 615 \\
$\Ga$ & 40 & 48{,}357 & 1{,}241{,}882 & 758 & 88{,}201{,}023 & 665 \\
$\Ga$ & 38 & 44{,}934 & 1{,}124{,}436 & 1{,}141 & 122{,}739{,}483 & 1{,}685 \\
$\Ga$ & 36 & 41{,}619 & 1{,}012{,}822 & 1{,}831 & 189{,}001{,}378 & 1{,}246 \\
$\Ga$ & 34 & 38{,}412 & 907{,}040 & 2{,}570 & 315{,}072{,}154 & 1{,}302 \\
$\Ga$ & 32 & 35{,}313 & 807{,}090 & 3{,}007 & 359{,}772{,}266 & 1{,}595 \\
$\Ga$ & 30 & 32{,}322 & 712{,}972 & 2{,}597 & 451{,}748{,}771 & 4{,}210 \\
$\Ga$ & 4 & 3{,}111 & 20{,}045 & 1{,}843 & 434{,}204{,}743 & 6{,}979 \\
$\Ga$ & 2 & 1{,}724 & 7{,}615 & 3 & 1{,}247{,}206 & 5 \\
\midrule
$\Gb$ & 50 & 67{,}092 & 1{,}916{,}592 & 73 & 17{,}037{,}541 & 129 \\
$\Gb$ & 48 & 63{,}129 & 1{,}769{,}986 & 460 & 42{,}390{,}272 & 687 \\
$\Gb$ & 46 & 59{,}274 & 1{,}629{,}212 & 1{,}991 & 120{,}499{,}030 & 599 \\
$\Gb$ & 44 & 55{,}527 & 1{,}494{,}270 & 3{,}476 & 363{,}213{,}799 & 8{,}930 \\
\midrule
$\Gsix$ & 66 & 116{,}540 & 4{,}279{,}152 & 218 & 59{,}635{,}574 & 349 \\
$\Gsix$ & 64 & 111{,}297 & 4{,}028{,}562 & 878 & 80{,}324{,}045 & 1{,}893 \\
\bottomrule
\end{tabular}
\caption{Target orders $k$ refuted with DRAT proofs accepted by \texttt{drat-trim}. The formula of order $k$ is unsatisfiable: no loopless cubic graph on $k$ vertices colors the graph with a vertex map of kind (O), (E0) or (E1). Times were measured while other computations shared the machine.}\label{tab:orders}
\end{table}


In kind (O) with $k$ close to 52 almost all fibers are singletons and the refutation is short; the
instances become harder as $k$ decreases and the fibers grow, and the certification of the final
formula dominates the running time. No instance required a connectivity or bridge clause: every
formula was already unsatisfiable without them, so the statement holds for colorings of kinds (O),
(E0), (E1) by arbitrary loopless cubic graphs, connected or not.

\section{Remarks and questions}\label{sec:remarks}

\emph{Smallest counterexamples.} A smallest counterexample to the Petersen coloring conjecture
belongs to $\Hthree$, and its order lies between 40 and 52~\cite{GJMMMU2026}. Membership of $\Ga$
and $\Gb$ in $\Hthree$ is necessary for them to be smallest counterexamples; it is not sufficient,
since $\Hthree$ may contain graphs of different orders that do not color one another.

\emph{The other known counterexamples.} The same question is open for the counterexample on 68
vertices and for the two on 112 vertices. An affirmative answer to ``is it colored by a smaller
bridgeless cubic graph'' would exhibit a smaller counterexample, possibly one of $\Ga$, $\Gb$. The
formulas of Section~\ref{sec:encoding} apply without change; the largest has target order 110.

\emph{Why the fiber lemmas matter.} Without Lemmas~\ref{lem:parity} and~\ref{lem:unused} the
unknown target has vertices and edges that no edge of $G$ constrains. For $\Ga$ the loop of
Section~\ref{sec:encoding} then produced, for each of the ten target orders from 32 to 50, about
two thousand targets with a bridge in half an hour, each excluded by one clause, without deciding
any order. With the lemmas every edge of the target is the image of an edge of $G$, and every
order was decided without a single such clause.

\emph{A question.} Is there a counterexample to the Petersen coloring conjecture that is colored
by a smaller counterexample? By Lemma~\ref{lem:parity} the fibers of such a coloring are all odd
or all even, and by Proposition~\ref{prop:nminus2} a coloring with every target vertex used needs
target order at most $n - 4$ when the graph has no triangles.

\section*{Data availability}

{\raggedright
The edge lists, encoders, checkers, runners, results and the hashes of all formulas and proofs
are in the repository \url{https://github.com/fsantibanezleal/CAOS_RESEARCH}, under
\path{problems/combinatorics/petersen-coloring/}. The computations of this paper are archived in
the folder \path{experiments/EXP-007-colorable-only-by-itself} of that directory (runner
\path{run_inc.py}, encoder \path{code/pcclib/hcolor.py}); the controls of
Section~\ref{sec:encoding} are in its \path{artifacts/agreement} folder. The solvers are CaDiCaL
1.7.3 and the CaDiCaL 1.9.5 build of PySAT; proofs were checked with \texttt{drat-trim}.
\par}

\begin{thebibliography}{99}
\bibitem{Jaeger1988} F.~Jaeger, Nowhere-zero flow problems, in: L.~W. Beineke, R.~J. Wilson
(eds.), Selected Topics in Graph Theory 3, Academic Press, London, 1988, 71--95.
\bibitem{Putman2026} B.~Putman, A 112-vertex counterexample to the Petersen Coloring Conjecture,
Zenodo v1.1.0 (2026), \href{https://doi.org/10.5281/zenodo.21845291}{doi:10.5281/zenodo.21845291};
arXiv:2608.10012.
\bibitem{GJMMMU2026} J.~Goedgebeur, J.~Jooken, E.~M\'a\v{c}ajov\'a, D.~Mattiolo, G.~Mazzuoccolo,
S.~Ulyanov, Disproving the Petersen Coloring Conjecture: theoretical analysis and an infinite
family of counterexamples, arXiv:2608.10028v3 (2026).
\bibitem{MMSW2025} Y.~Ma, D.~Mattiolo, E.~Steffen, I.~H. Wolf, Sets of $r$-graphs that color all
$r$-graphs, Combinatorica 45 (2025), Article 16,
\href{https://doi.org/10.1007/s00493-025-00144-4}{doi:10.1007/s00493-025-00144-4}; arXiv:2305.08619.
\bibitem{Fleischner1992} H.~Fleischner, Spanning eulerian subgraphs, the splitting lemma, and
Petersen's theorem, Discrete Math. 101 (1992), 33--37.
\bibitem{KaiserKLW2007} T.~Kaiser, R.~Ku\v{z}el, H.~Li, G.~Wang, A note on $k$-walks in bridgeless
graphs, Graphs Combin. 23 (2007), 303--308.
\bibitem{HoG} K.~Coolsaet, S.~D'hondt, J.~Goedgebeur, House of Graphs 2.0: a database of
interesting graphs and more, Discrete Appl. Math. 325 (2023), 97--107;
\url{https://houseofgraphs.org}.
\bibitem{CaDiCaL} A.~Biere, T.~Faller, K.~Fazekas, M.~Fleury, N.~Froleyks, F.~Pollitt, CaDiCaL 2.0,
in: Computer Aided Verification (CAV 2024), Lecture Notes in Comput. Sci. 14681, Springer, 2024,
133--152.
\bibitem{DRATtrim} N.~Wetzler, M.~J.~H. Heule, W.~A. Hunt Jr., DRAT-trim: efficient checking and
trimming using expressive clausal proofs, in: Theory and Applications of Satisfiability Testing
(SAT 2014), Lecture Notes in Comput. Sci. 8561, Springer, 2014, 422--429.
\end{thebibliography}


\end{document}
